\section{Learning Finite Joint Recovery}
\label{sec:method}

The immediate term $q_a$ is known once an action is proposed; learning it is
unnecessary. The unknown quantity is the reward of mutually compatible
replacements returned by a finite execution procedure, conditional on its
scope, native work allowance, and actual warm start. We predict this joint
response and use it to prioritize requests. Scope construction, clique
inequalities, and cost-aware destroy--repair allocation supply the execution
framework; the learned component estimates finite joint recovery within it.

\subsection{A Shared Conditional Response Encoder}
\label{sec:v4-response}
For action $a=(C_a,E_a)$, the fixed base is
$B_a=(S\setminus((N(C_a)\cap S)\cup E_a))\cup C_a$.
The eligible set and nested prefixes $R_{a,r}$ follow
Eq.~\eqref{eq:v4-recovery-scope}. Commitments use a common pool of the top
32 vertices by individual immediate gain and 32 seeded exploratory vertices;
a compatible prefix supplies $C_a$. Optional releases follow a seeded
breadth-first traversal of the incumbent boundary graph, in which two
incumbent vertices are adjacent when they share an unselected neighbor.
These proposals use no recovery outcome. Identical actual prefixes retain
their lowest-cap representative, with no repeated alias queries.

Write $R=R_{a,r}$ and let $W\subseteq R$ be the actual feasible warm start.
We construct the induced conflict graph $G[R]$ and factors $\mathcal F_R$:
greedily grown conflict cliques in decreasing reward--degree order with
identifier ties, supplied resource cliques restricted to $R$, and size-two
factors for uncovered edges. All factors are genuine conflict cliques and
duplicates are removed. There is no maximal-clique enumeration. Let
$s_G=\max(1,|V|^{-1}\sum_{v\in V}w_v)$ and $d_i$ be the degree in $G[R]$.
Node features contain $w_i/s_G$, $\log(1+d_i)$,
$d_i/\max(1,|R|-1)$, $\mathbf1[i\in W]$, and the logarithms of one plus
the incident-factor count and maximum incident-factor size (at least one).

The encoder has width 32 and two layers. Starting from an affine embedding
and ReLU, each layer combines a self transformation, mean conflict-neighbor
messages, and a mean relay through incident factors:
\begin{equation}
 h_i^{\ell+1}=\operatorname{ReLU}\!\left(
 A_\ell h_i^\ell+
 B_\ell\operatorname{mean}_{j\in N_R(i)}h_j^\ell+
 D_\ell\operatorname{mean}_{F\ni i}\bar h_F^\ell\right),
 \label{eq:v4-encoder}
\end{equation}
where $\bar h_F^\ell=|F|^{-1}\sum_{j\in F}h_j^\ell$ and empty means are
zero. The structural unit is $(R,\mathcal F_R,W,s_G)$: requests with the
same unit share one encoding within a decision. Neither $q_a$, native work
allowance, remaining time, action identifier, nor owner identifier enters
this encoder. A changed actual warm set requires a new encoding. This
factorization follows the recovery interface, whose search input is the
weighted induced graph and warm membership; the base already determines
eligibility.

For workpoint $\kappa=(k,b)$ and remaining time $\tau$, the request head uses
\begin{equation}
 c_{\kappa,\tau,W}=\left(
 \mathbf1[k=\mathrm{seconds}],\mathbf1[k=\mathrm{nodes}],
 \mathbf1[k=\mathrm{iterations}],\log(1+b),\log(1+\tau),
 \frac{w(W)}{s_G},\frac{|W|}{\max(1,|R|)}\right).
 \label{eq:v4-head-context}
\end{equation}
The head also receives observable scope summaries: normalized reward
sum/maximum/mean/deviation, logarithmic edge and factor counts, edge density,
and factor size
statistics, mean memberships, the fraction of factors of size at least
three, a uniform-occupancy capacity summary, and a disjoint-clique upper
bound. Logarithmic counts use $\log(1+\cdot)$. The calibrated
request cost and current best gain are allocation quantities, not head
inputs. A shared width-32 ReLU head applied to
$(h_i^2,\operatorname{mean}_{j\in R}h_j^2,c_{\kappa,\tau,W},s_R)$ produces
the occupancy logit $\ell_i$.

\subsection{A Capacity-Informed Value Readout}
\label{sec:v4-capacity}
Set $p_i=\sigma(\ell_i)$. Using loads computed once from these proposed
occupancies, apply one simultaneous contraction:
\begin{equation}
 L_F=\sum_{j\in F}p_j,\qquad
 z_i=\frac{p_i}{\max\{1,\max_{F\ni i}L_F\}}.
 \label{eq:v4-capacity}
\end{equation}
The inner maximum is zero without incident factors. Every factor then has
load at most one: if $L_F>1$, each of its node denominators is at least
$L_F$; otherwise contraction cannot increase load. Edge coverage gives
$z_i+z_j\leq1$ on every conflict edge. This uses classical clique
capacities~\cite{chvatal1975polytopes} as an inductive bias for a value
surrogate. It is not an integer recovery or a certificate of its value:
half-occupancy on a unit-weight five-cycle satisfies the factors but sums
to $2.5$, whereas an independent set has weight at most two.

The signed normalized prediction and deployed native score are
\begin{equation}
 \widehat t_u=\frac{q_a}{s_G}+\sum_{i\in R}\frac{w_i}{s_G}z_i,
 \qquad \widehat g_u=s_G[\widehat t_u]_+,
 \label{eq:v4-prediction}
\end{equation}
for request $u=(a,r,\kappa,W,\tau)$. Thus $q_a$ appears only as the known
final offset; training uses $\widehat t_u$ before clamping. An empty pool
has zero recovery contribution. Joint incompatibilities affect both the
representations and the aggregate occupancy, rather than allowing each
individually attractive replacement to contribute its full reward.

\subsection{Supervision from Actual Finite Responses}
\label{sec:v4-supervision}
Alternatives are executed independently from the same immutable state,
with the same executor reset and their declared workpoints. Each alternative
restores that state's remaining allowance and actual warm membership;
its outcome never warms another alternative. The primary executor is
CHILS~\cite{grossmann2025chils} with population one, one thread, seed 17,
native warm input, and a reset for each request. Two snapshots each receive
weight $1/2$: after paid validation of known $S\cap R$ recoveries, and
after at most one actually startable round-robin request at cap 64 and
50\,ms under the same history deadline.

A valid returned membership $T_u$ supplies the signed target
$t_u=(q_a+w(T_u))/s_G$ and binary labels $\mathbf1[i\in T_u]$. Its
admitted target $y_u$ is $[t_u]_+$ when full validation and rescoring finish
within the cloned remaining allowance, and zero otherwise. Valid late or
nonimproving outputs still supervise signed values and masks. Failed outputs
have no invented value or mask target, but remain in admitted-outcome
ranking with $y_u=0$.

Let $b_s$ be the snapshot's best materialized gain divided by $s_G$, and
let $\mathcal P_s=\{(u,v):[y_u-b_s]_+>[y_v-b_s]_+\}$. The state loss is
\begin{align}
 \mathcal L_s={}&\operatorname{mean}_{u\in\mathcal V_s}
       \operatorname{SmoothL1}(\widehat t_u,t_u)\nonumber\\
 &+\operatorname{mean}_{(u,v)\in\mathcal P_s}
       \log\!\left(1+e^{-(\widehat t_u-\widehat t_v)/0.2}\right)
   +0.05\mathcal L_{\mathrm{mask}},
 \label{eq:v4-loss}
\end{align}
where $\mathcal V_s$ contains valid returned outputs. Mask BCE is averaged
over nodes within each valid nonempty request and then over requests.
Empty terms are zero. Both value regression and pair differences use signed
pre-clamp predictions. Domains have equal total weight, graphs have equal
weight within a domain, and the two snapshot weights remain fixed; failed
graph slots and unavailable snapshots contribute zero without redistributing
their weight. The protocol uses Adam at $10^{-3}$ for 40 epochs in batches
of four graphs and fit seeds 17, 29, and 43. Checkpoints use the first minimum
validation conditional allocation regret: this offline, single-request
proxy can abstain and does not measure the online multi-call controller.

Matched controls replace the capacity readout with a free scalar head,
remove only mask BCE, mask all warm-input channels, or use the same
observable summaries without graph messages. These isolate readout,
supervision, warm conditioning, and representation effects.

\subsection{Paid Request Allocation and Feasible Admission}
\label{sec:v4-controller}
All proposals, scope/factor preparation, warm validation, transfers,
encoding, cache access, heads, solver calls, and final checks are charged
to the decision. Native slices are 10, 50, and 200\,ms; their common
calibrated full request costs are 35, 75, and 214\,ms.
Calibration is not a runtime guarantee. Only unspent requests whose native
slice and calibrated cost fit the current remaining time are startable.
With best materialized gain $g^*$, the learned policy selects
\begin{equation}
 u^*=\arg\max_u\left(
 \frac{[\widehat g_u-g^*]_+}{\widehat c_u},
 \frac{\widehat g_u}{\widehat c_u},-\operatorname{index}(u)\right)
 \label{eq:v4-priority}
\end{equation}
lexicographically. This point-value heuristic is neither an optimal
allocation nor an estimate of expected positive marginal gain. Zero first
components do not cause online abstention.

\begin{algorithm}[t]
\caption{Finite joint-recovery allocation}
\label{alg:v4-recovery}
\begin{algorithmic}[1]
\State Start deadline; validate $S$; set $S^*\gets S$, $g^*\gets0$
\State Prepare fixed actions, unique nested scopes, and requests
\State Validate known $S\cap R$; retain timely improving schedules
\State Initialize same-action warm memory from timely known recoveries
\While{fewer than eight calls and a startable unspent request exists}
  \State Score current request views; choose by Eq.~\eqref{eq:v4-priority}
  \State Recheck time; mark chosen request spent before launch
  \State Execute with actual same-action warm memory intersected with $R$
  \State Fully validate $T\subseteq R$ and rescore $B_a\cup T$
  \If{$T$ is valid and validation finishes before the deadline}
    \State Retain better recovery warm memory and improving $S^*$
  \EndIf
\EndWhile
\State Fully validate and return $S^*$; measure caller receipt time
\end{algorithmic}
\end{algorithm}

Every launched request is spent, including invalid, failed, and late calls.
Only timely validated recoveries can update same-action warm memory or the
best feasible schedule; all policies preserve their actual incumbent.
The primary gain uses the fully validated schedule actually received by
the caller strictly before $D$, reverting to the original $S$ otherwise.
A separately observed schedule validated inside the procedure before $D$
is a secondary availability measure, not a return-time guarantee.
Soft solver limits therefore require reporting actual latency and misses.
